% problem_id: S046-thm:equiv-iso-adj
% source_id: S046 (HoTT Book)
% source_file: equivalences.tex
% statement_locator: equivalences.tex lines 215--217
% proof_locator: equivalences.tex lines 218--242
% env: thm
% pairing: tight (verified)
% status: complete (statement + author proof verbatim + reason + locators)
%
\begin{thm}\label{thm:equiv-iso-adj}
  For any $f:A\to B$ we have $\qinv(f)\to\ishae(f)$.
\end{thm}

\begin{proof}
Suppose that $(g,\eta,\epsilon)$ is a quasi-inverse for $f$. We have to provide
a quadruple $(g',\eta',\epsilon',\tau)$ witnessing that $f$ is a half adjoint equivalence. To
define $g'$ and $\eta'$, we can just make the obvious choice by setting $g'
\defeq g$ and $\eta'\defeq \eta$. However, in the definition of $\epsilon'$ we
need start worrying about the construction of $\tau$, so we cannot just follow our nose
and take $\epsilon'$ to be $\epsilon$. Instead, we take
\begin{equation*}
\epsilon'(b) \defeq \opp{\epsilon(f(g(b)))}\ct (\ap{f}{\eta(g(b))}\ct \epsilon(b)).
\end{equation*}
Now we need to find
\begin{equation*}
\tau(a): \ap{f}{\eta(a)}=\opp{\epsilon(f(g(f(a))))}\ct (\ap{f}{\eta(g(f(a)))}\ct \epsilon(f(a))).
\end{equation*}
Note first that by \cref{cor:hom-fg}, we have
%$\eta(g(f(a)))\ct\eta(a)=\ap{g}{\ap{f}{\eta(a)}}\ct\eta(a)$ and hence it follows that
$\eta(g(f(a)))=\ap{g}{\ap{f}{\eta(a)}}$. Therefore, we can apply
\cref{lem:htpy-natural} to compute
\begin{align*}
\ap{f}{\eta(g(f(a)))}\ct \epsilon(f(a))
& = \ap{f}{\ap{g}{\ap{f}{\eta(a)}}}\ct \epsilon(f(a))\\
& = \epsilon(f(g(f(a))))\ct \ap{f}{\eta(a)}
\end{align*}
from which we get the desired path $\tau(a)$.
\end{proof}
